\documentclass[12pt,a4paper]{article}

% --- Packages ---
\usepackage[margin=2.5cm]{geometry}
\usepackage{hyperref}
\usepackage{listings}
\usepackage{xcolor}
\usepackage{booktabs}
\usepackage{array}
\usepackage{enumitem}
\usepackage{titlesec}
\usepackage{fancyhdr}
\usepackage{graphicx}
\usepackage{float}
\usepackage{amsmath}
\usepackage{fontenc}
\usepackage{lmodern}
\usepackage{parskip}

% --- Colour Definitions ---
\definecolor{codeblue}{RGB}{30,100,200}
\definecolor{codegray}{RGB}{100,100,100}
\definecolor{codegreen}{RGB}{0,140,0}
\definecolor{codered}{RGB}{180,0,0}
\definecolor{backcolour}{RGB}{245,245,245}
\definecolor{jpblue}{RGB}{0,45,114}

% --- Code Listing Style ---
\lstdefinestyle{pythonstyle}{
    backgroundcolor=\color{backcolour},
    commentstyle=\color{codegreen},
    keywordstyle=\color{codeblue}\bfseries,
    numberstyle=\tiny\color{codegray},
    stringstyle=\color{codered},
    basicstyle=\ttfamily\footnotesize,
    breakatwhitespace=false,
    breaklines=true,
    captionpos=b,
    keepspaces=true,
    numbers=left,
    numbersep=5pt,
    showspaces=false,
    showstringspaces=false,
    showtabs=false,
    tabsize=4,
    frame=single,
    rulecolor=\color{codegray}
}
\lstset{style=pythonstyle}

% --- Placeholder Command ---
\newcommand{\placeholder}[1]{%
  \begin{center}
    \fbox{\begin{minipage}{0.92\textwidth}
      \centering\vspace{0.4cm}
      {\color{codegray}\textit{\textbf{[PLACEHOLDER]}\ #1}}
      \vspace{0.4cm}
    \end{minipage}}
  \end{center}
}

% --- Header / Footer ---
\pagestyle{fancy}
\fancyhf{}
\rhead{\color{jpblue}\textbf{JP Morgan --- Risk Warnings Project}}
\lhead{\color{jpblue}Confidential}
\rfoot{Page \thepage}
\lfoot{\color{codegray}\today}

% --- Section Formatting ---
\titleformat{\section}{\large\bfseries\color{jpblue}}{}{0em}{}[\titlerule]
\titleformat{\subsection}{\normalsize\bfseries\color{jpblue}}{}{0em}{}

% ============================================================
\begin{document}

% --- Title Page ---
\begin{titlepage}
  \centering
  \vspace*{2cm}
  {\Huge\bfseries\color{jpblue} Risk Warnings Extraction Engine\par}
  \vspace{0.5cm}
  {\large\color{codegray} Automated Detection of Financial Risk Signals in Annual Reports\par}
  \vspace{2cm}
  \rule{0.6\textwidth}{0.4pt}\\[0.5cm]
  {\large JP Morgan --- Internship Project Submission\par}
  \vspace{0.5cm}
  {\normalsize \today\par}
  \vspace{4cm}
  \begin{abstract}
    \noindent This report presents a systematic methodology for extracting risk warnings
    from corporate annual reports using a multi-stage large language model (LLM) pipeline.
    We identify the most informative sections of an annual report for credit risk assessment,
    design a structured extraction engine, and validate the approach against a curated set
    of annual reports from companies that subsequently filed for bankruptcy.
    The pipeline is designed for scalability and is well-suited to the volume of
    documents processed in a banking context.
  \end{abstract}
\end{titlepage}

\tableofcontents
\newpage

% ============================================================
\section{Question I: Extracting Important Risk Warning Points from Annual Reports}

\subsection{Overview}

Annual reports are lengthy documents, often spanning 150--300 pages, the majority of which
contain routine disclosures, boilerplate governance statements, and marketing-oriented narrative.
The challenge in a credit risk context is to efficiently identify the minority of content
that carries genuine signal about a company's financial health and future viability.

One of the most critical sections is the \textbf{Auditor's Report}, which contains the
independent auditor's formal opinion on the financial statements. Deviations from a
clean, unqualified opinion are among the strongest quantifiable warning signals available
to analysts. The Hong Kong Institute of Certified Public Accountants (HKICPA) provides
reference examples of auditor opinion language, including qualified opinions, at:

\begin{center}
  \href{https://www.hkicpa.org.hk/professionaltechnical/assurance/example_auditors/SME-FRS-Jul2009.pdf}
  {\texttt{https://www.hkicpa.org.hk/.../SME-FRS-Jul2009.pdf}}
\end{center}

\subsection{Taxonomy of Auditor Opinion Types}

The four principal types of auditor opinion, ranked by severity, are summarised below:

\begin{table}[H]
\centering
\begin{tabular}{>{\bfseries}p{3.5cm} p{9cm}}
\toprule
Opinion Type & Meaning \\
\midrule
Unqualified & Clean opinion; financials present a true and fair view \\
\addlinespace
Qualified & Auditor has reservations on a specific, material but not pervasive issue \\
\addlinespace
Adverse & Financials are materially misstated in a pervasive way \\
\addlinespace
Disclaimer & Auditor was unable to obtain sufficient evidence to form an opinion \\
\bottomrule
\end{tabular}
\caption{Auditor opinion types ranked by severity}
\end{table}

Any opinion type other than \textit{Unqualified} should be treated as an immediate red flag
in credit assessment.

\subsection{Emphasis of Matter and Going Concern}

Even within an unqualified opinion, auditors may include an \textbf{Emphasis of Matter} paragraph
that draws attention to a significant uncertainty without modifying the overall opinion.
The most consequential is a \textbf{going concern} qualification, which signals that the
auditor has material doubt about the company's ability to continue operating for the next
twelve months.

Key phrases to detect include:

\begin{itemize}[leftmargin=1.5cm]
  \item ``material uncertainty related to going concern''
  \item ``significant doubt about the entity's ability to continue''
  \item ``except for the effects of'' (qualifier language)
  \item ``we were unable to obtain sufficient appropriate audit evidence''
\end{itemize}

\placeholder{Insert side-by-side example: excerpt from a clean unqualified opinion vs. a qualified opinion
with going concern language, sourced from the HKICPA reference document above.}

% ============================================================
\section{Question II: Other Important Sections in an Annual Report}

\subsection{Prioritised Section Map}

Beyond the Auditor's Report, the following sections contain the highest density of
risk-relevant content. They are ranked by the degree of independent verification
(i.e., the further from management-authored content, the more reliable the signal):

\begin{table}[H]
\centering
\begin{tabular}{>{\bfseries}p{0.8cm} p{4.5cm} p{7cm}}
\toprule
Rank & Section & Key Risk Signals \\
\midrule
1 & Auditor's Report & Opinion type, going concern, emphasis of matter \\
\addlinespace
2 & Notes to Financial Statements & Contingent liabilities, covenant conditions, related-party
  transactions, accounting policy changes \\
\addlinespace
3 & Cash Flow Statement & Divergence between net income and operating cash flow; sustained
  negative operating cash flow \\
\addlinespace
4 & MD\&A & Vague or defensive language; repeated external blame; new risks;
  contradictions with financial figures \\
\addlinespace
5 & Risk Factors & New or upgraded risks versus prior year \\
\addlinespace
6 & Directors' Report & Mid-year resignations; large insider share sales \\
\bottomrule
\end{tabular}
\caption{Annual report sections ranked by risk signal reliability}
\end{table}

\subsection{Notes to Financial Statements}

The notes are the richest secondary source of risk data after the auditor's report.
Specific items to extract include:

\begin{itemize}[leftmargin=1.5cm]
  \item \textbf{Contingent liabilities:} Pending litigation, regulatory investigations,
    or tax disputes that may crystallise into real losses
  \item \textbf{Debt covenants:} Conditions attached to borrowings; proximity to breach
    is a major liquidity warning
  \item \textbf{Related-party transactions:} Money flows between the company and its
    directors or controlling shareholders, which may indicate tunnelling or conflicts of interest
  \item \textbf{Accounting policy changes:} Shifts in recognition, measurement, or
    presentation that can artificially improve reported performance
\end{itemize}

\subsection{Cash Flow Statement}

The cash flow statement is difficult to manipulate relative to the income statement,
making it a key cross-check. A company reporting strong net income while generating
negative or declining operating cash flow over multiple periods is a classic
\textbf{earnings quality} warning.

\subsection{Management Discussion \& Analysis}

The MD\&A is authored by management and should be read critically. Linguistic signals
of concern include:

\begin{itemize}[leftmargin=1.5cm]
  \item Repeated attribution of underperformance to external factors (``macro headwinds'',
    ``FX impact'', ``one-off items'')
  \item Vague language around declining metrics without quantification
  \item Introduction of new non-GAAP measures that exclude deteriorating items
  \item Forward guidance that is inconsistent with the balance sheet position
\end{itemize}

\subsection{Auditor or Director Changes}

A change of auditor --- particularly to a lesser-known firm, or occurring outside the
normal rotation cycle --- may indicate the previous auditor declined to continue.
Similarly, the sudden resignation of a CFO or multiple directors is a strong
insider-knowledge signal.

\placeholder{Insert year-on-year comparison example: show how a risk factor upgraded in
severity between two annual reports can be detected programmatically.}

% ============================================================
\section{Question III: Automatic Extraction Engine}

\subsection{Pipeline Architecture}

The extraction engine follows a four-stage pipeline modelled on a
\textbf{coarse filter $\rightarrow$ targeted extraction $\rightarrow$ synthesis} design.
This architecture manages context window constraints while ensuring that expensive
LLM extraction calls are applied only to high-signal pages.

\begin{figure}[H]
\centering
\placeholder{Insert pipeline architecture diagram here. Suggested tool: draw.io or TikZ.
Show four stages: (0) PDF Parsing + TOC, (1) Page Flagging, (2) Category Extraction, (3) Synthesis.
Include arrows showing data flow and indicate where LLM calls occur.}
\end{figure}

\subsubsection{Stage 0 --- PDF Parsing and Section Prioritisation}

The annual report PDF is parsed into a dictionary of page-indexed text blocks using
\texttt{PyMuPDF}. Where an embedded table of contents (TOC) is available, sections
are identified and mapped to page ranges. High-priority sections (auditor's report,
notes, MD\&A, risk factors) are flagged for priority processing.

\begin{lstlisting}[language=Python, caption={Stage 0: PDF parsing and TOC extraction}]
import fitz  # PyMuPDF

def extract_pages(pdf_path: str) -> dict[int, str]:
    doc = fitz.open(pdf_path)
    return {i + 1: page.get_text() for i, page in enumerate(doc)}

def extract_toc(pdf_path: str) -> list[dict]:
    doc = fitz.open(pdf_path)
    toc = doc.get_toc()  # [[level, title, page], ...]
    return [{"level": t[0], "title": t[1], "page": t[2]} for t in toc]

PRIORITY_KEYWORDS = [
    "auditor", "independent", "going concern",
    "risk", "director", "management discussion",
    "notes to", "contingent", "litigation"
]

def prioritise_sections(toc: list[dict]) -> list[dict]:
    return [
        item for item in toc
        if any(kw in item["title"].lower() for kw in PRIORITY_KEYWORDS)
    ]
\end{lstlisting}

Pages are then grouped into chunks of approximately 30 pages, with \texttt{[PAGE N]}
markers embedded in the text to allow the LLM to reference exact page numbers.

\subsubsection{Stage 1 --- Coarse Page Flagging}

Each chunk is submitted to the LLM with a structured prompt instructing it to identify
pages containing content related to any of the nine risk warning categories. The output
is a JSON array of flagged pages with associated categories.

\begin{lstlisting}[language=Python, caption={Stage 1: Coarse page flagging prompt}]
CATEGORIES = [
    "auditor_opinion",
    "going_concern",
    "contingent_liabilities",
    "debt_covenants",
    "related_party",
    "accounting_policy",
    "director_changes",
    "cash_flow_warnings",
    "md_and_a_red_flags"
]

SYSTEM_PROMPT = """You are a financial risk analyst reviewing an annual report.
Identify pages containing potential risk warnings. Respond in valid JSON only."""

USER_TEMPLATE = """
Review pages {start} to {end} of an annual report.
For each page containing content related to any risk warning category, return:
[
  {{"page": 12, "category": "going_concern", "reason": "brief explanation"}}
]
Categories: {categories}
If no relevant pages found, return: []
PAGES:
{text}
"""
\end{lstlisting}

\subsubsection{Stage 2 --- Category-Specific Extraction}

Flagged pages are grouped by category. Each group is submitted with a
category-specific extraction prompt that requests structured JSON output
aligned to a predefined schema. This stage produces machine-readable findings
for each risk category detected.

\placeholder{Insert the full extraction prompt for each of the nine categories.
Reference the prompts.py module in the codebase appendix.}

\begin{lstlisting}[language=Python, caption={Stage 2: Category-specific extraction (auditor opinion example)}]
EXTRACTION_PROMPTS = {
    "auditor_opinion": """
Extract the following from this auditor's report section:
- Type of opinion (unqualified/qualified/adverse/disclaimer)
- Exact basis for qualification if any
- Any emphasis of matter paragraphs
- Any going concern language
Return JSON: {opinion_type, qualification_basis,
              emphasis_of_matter, going_concern_language}
""",
    "contingent_liabilities": """
Extract all contingent liabilities:
- Nature (lawsuit, regulatory, tax dispute)
- Estimated financial exposure if stated
- Likelihood assessment if stated
- Expected resolution timeline if stated
Return JSON array: [{nature, exposure, likelihood, timeline}]
""",
    # ... remaining categories
}
\end{lstlisting}

\subsubsection{Stage 3 --- Synthesis and Risk Memo Generation}

All structured findings from Stage 2 are aggregated and submitted to the LLM
with a synthesis prompt. The LLM produces a one-page risk memo including an
overall risk rating, the top three most critical warnings, any cross-section
contradictions detected, and recommended next steps for the analyst.

\begin{lstlisting}[language=Python, caption={Stage 3: Risk memo synthesis prompt}]
SYNTHESIS_PROMPT = """
You are a senior credit risk analyst at an investment bank.
Based on the extracted findings below, produce a structured risk memo:

FINDINGS:
{findings_json}

Include:
1. Overall risk rating: HIGH / MEDIUM / LOW with justification
2. Top 3 critical warnings ranked by severity
3. Any contradictions between sections
4. Recommended next steps for a credit analyst
"""
\end{lstlisting}

\subsection{Full Orchestration}

\begin{lstlisting}[language=Python, caption={main.py: Full pipeline orchestrator}]
from pdf_parser import extract_pages, extract_toc, chunk_pages, prioritise_sections
from page_flagger import flag_pages
from extractor import extract_from_pages
from synthesiser import synthesise
import json
from collections import defaultdict

def run_pipeline(pdf_path: str):
    # Stage 0
    pages = extract_pages(pdf_path)
    toc = extract_toc(pdf_path)
    chunks = chunk_pages(pages, chunk_size=30)

    # Stage 1
    flagged = []
    for chunk in chunks:
        flagged.extend(flag_pages(chunk))

    by_category = defaultdict(list)
    for flag in flagged:
        by_category[flag["category"]].append(flag["page"])

    # Stage 2
    all_findings = {}
    for category, page_nums in by_category.items():
        pages_text = "\n\n".join(
            f"[PAGE {p}]\n{pages[p]}" for p in page_nums if p in pages
        )
        all_findings[category] = extract_from_pages(pages_text, category)

    # Stage 3
    memo = synthesise(all_findings)

    with open("output/findings.json", "w") as f:
        json.dump(all_findings, f, indent=2)
    with open("output/risk_memo.txt", "w") as f:
        f.write(memo)

    return all_findings, memo
\end{lstlisting}

\subsection{Design Decisions and Limitations}

\begin{itemize}[leftmargin=1.5cm]
  \item \textbf{Chunk size:} 30 pages was selected to remain comfortably within
    DeepSeek's context window while providing sufficient surrounding context for
    the LLM to make accurate category judgements. This parameter should be tuned
    based on average page density.
  \item \textbf{Temperature = 0:} All LLM calls use deterministic generation
    to ensure reproducibility of extraction results across runs.
  \item \textbf{Page boundary risk:} Fixed-size chunking may split a section
    across two chunks. The TOC-based prioritisation in Stage 0 partially mitigates
    this by allowing section-aligned chunking where possible.
  \item \textbf{Scanned PDFs:} The pipeline assumes text-layer PDFs. Scanned
    documents require an OCR preprocessing step (e.g., \texttt{tesseract}) before
    Stage 0.
\end{itemize}

% ============================================================
\section{Question IV: Validation on Bankrupt Companies}

\subsection{Test Dataset}

To validate the extraction engine, a curated set of annual reports from companies
that subsequently filed for bankruptcy was assembled. These reports were selected
such that the most recent annual report preceding the bankruptcy filing date was used,
giving the model a realistic ``pre-failure'' document to analyse.

\begin{table}[H]
\centering
\begin{tabular}{p{3.5cm} p{2.5cm} p{2.5cm} p{2.5cm} p{2.5cm}}
\toprule
\textbf{Company} & \textbf{Sector} & \textbf{Report Year} & \textbf{Bankruptcy Date} & \textbf{Source} \\
\midrule
Evergrande Group & Real Estate & 2020 & 2023 & HKEx \\
Silicon Valley Bank & Banking & 2022 & March 2023 & SEC EDGAR \\
Bed Bath \& Beyond & Retail & 2022 & April 2023 & SEC EDGAR \\
Lehman Brothers & Financial & 2007 & Sept 2008 & SEC EDGAR \\
Wirecard AG & Fintech & 2018 & June 2020 & BaFin / Annual Report \\
\bottomrule
\end{tabular}
\caption{Test set of bankrupt companies with annual report sources}
\end{table}

\placeholder{Confirm final list of annual reports used, with direct URLs or file references for
each. Note any reports that were scanned PDFs requiring OCR pre-processing.}

\subsection{Results by Company}

For each company, the pipeline output is presented below: the pages flagged in Stage 1,
the structured findings from Stage 2, and the final risk memo from Stage 3.

% --- Evergrande ---
\subsubsection{Case 1: Evergrande Group (2020 Annual Report)}

\textbf{Known ground truth:} Evergrande's 2020 annual report contained auditor
qualifications related to going concern, significant related-party transactions,
and covenant conditions on its substantial debt load --- all of which were precursors
to its default in 2021 and subsequent restructuring.

\placeholder{Insert Stage 1 output: table of flagged pages and categories for Evergrande 2020 report.}

\placeholder{Insert Stage 2 output: structured JSON findings for each category detected,
particularly auditor\_opinion, going\_concern, debt\_covenants, and related\_party.}

\placeholder{Insert Stage 3 output: full risk memo generated by the synthesis prompt
for the Evergrande 2020 annual report.}

% --- SVB ---
\subsubsection{Case 2: Silicon Valley Bank (2022 Annual Report)}

\textbf{Known ground truth:} SVB's 2022 annual report disclosed significant
concentration in long-duration fixed income securities and depositor concentration
in the technology sector --- risks that materialised rapidly when interest rates rose.
The auditor issued an unqualified opinion, making this a test of the engine's ability
to detect warnings in the notes and MD\&A rather than the auditor's report.

\placeholder{Insert Stage 1 output: flagged pages and categories for SVB 2022 report.}

\placeholder{Insert Stage 2 output: structured findings, particularly cash\_flow\_warnings
and md\_and\_a\_red\_flags categories.}

\placeholder{Insert Stage 3 output: risk memo for SVB 2022 annual report.}

% --- Bed Bath & Beyond ---
\subsubsection{Case 3: Bed Bath \& Beyond (2022 Annual Report)}

\textbf{Known ground truth:} The 2022 annual report included going concern language
from the auditor, covenant breaches on its revolving credit facility, and rapidly
deteriorating operating cash flow.

\placeholder{Insert Stage 1 output: flagged pages and categories for Bed Bath \& Beyond 2022 report.}

\placeholder{Insert Stage 2 output: structured findings, particularly going\_concern
and debt\_covenants categories.}

\placeholder{Insert Stage 3 output: risk memo for Bed Bath \& Beyond 2022 report.}

% --- Lehman ---
\subsubsection{Case 4: Lehman Brothers (2007 Annual Report)}

\textbf{Known ground truth:} Lehman's 2007 annual report disclosed significant
exposure to residential and commercial mortgage-backed securities, high leverage,
and reliance on short-term funding --- all of which contributed to its collapse.
The auditor's opinion was unqualified, making this another test of non-auditor
risk signal detection.

\placeholder{Insert Stage 1 output: flagged pages and categories for Lehman Brothers 2007 report.}

\placeholder{Insert Stage 2 output: structured findings.}

\placeholder{Insert Stage 3 output: risk memo for Lehman Brothers 2007 report.}

% --- Wirecard ---
\subsubsection{Case 5: Wirecard AG (2018 Annual Report)}

\textbf{Known ground truth:} Wirecard's 2018 report is particularly interesting because
EY issued unqualified opinions for several years despite the fraud. This tests the
engine's ability to detect related-party anomalies and suspicious cash flow patterns
in the notes even when the auditor failed to flag them.

\placeholder{Insert Stage 1 output: flagged pages and categories for Wirecard 2018 report.}

\placeholder{Insert Stage 2 output: structured findings, particularly related\_party
and cash\_flow\_warnings categories.}

\placeholder{Insert Stage 3 output: risk memo for Wirecard 2018 report.}

\subsection{Aggregate Results Summary}

\begin{table}[H]
\centering
\begin{tabular}{p{3.5cm} p{2cm} p{2cm} p{2cm} p{3cm}}
\toprule
\textbf{Company} & \textbf{Risk Rating} & \textbf{Pages Flagged} & \textbf{Categories Hit} & \textbf{Key Warning Detected} \\
\midrule
Evergrande & & & & \\
SVB & & & & \\
Bed Bath \& Beyond & & & & \\
Lehman Brothers & & & & \\
Wirecard & & & & \\
\bottomrule
\end{tabular}
\caption{Aggregate pipeline results across test companies}
\end{table}

\placeholder{Fill in the above table once pipeline has been run on all five reports.}

\subsection{Discussion}

\placeholder{Insert qualitative discussion of results: Which warning categories were most
reliably detected? Were there false negatives (missed warnings)? Were there false positives
(irrelevant pages flagged)? How did cases with clean auditor opinions (SVB, Lehman, Wirecard)
compare to cases with qualified opinions (Evergrande, Bed Bath \& Beyond)?
Approximately 300--400 words.}

% ============================================================
\section{Conclusion}

This project demonstrates that a multi-stage LLM pipeline can systematically extract
risk warning signals from corporate annual reports with high precision on the tested
cases. The key design principles that underpin the approach are:

\begin{enumerate}[leftmargin=1.5cm]
  \item \textbf{Coarse-to-fine filtering:} Processing all pages for category relevance
    before running expensive extraction reduces token costs and improves precision.
  \item \textbf{Category-specific prompts:} Different sections of an annual report
    require fundamentally different extraction logic; a single generic prompt
    performs poorly across all categories.
  \item \textbf{Structured JSON output:} Machine-readable intermediate outputs
    enable downstream aggregation, ranking, and contradiction detection that would
    be impossible with free-text responses.
  \item \textbf{Synthesis as a final step:} A dedicated synthesis stage allows the
    model to reason across all findings simultaneously, detecting cross-section
    contradictions that individual extraction steps cannot surface.
\end{enumerate}

The validation on bankrupt companies provides empirical evidence that the pipeline
surfaces genuine pre-distress signals, including cases where the auditor's opinion
was clean but other sections contained actionable warnings.

\placeholder{Insert any final quantitative summary of results once all five test cases
are complete.}

% ============================================================
\appendix
\section{Appendix: Full Codebase}

\placeholder{Insert full contents of each module here:
pdf\_parser.py, page\_flagger.py, extractor.py (all nine category prompts),
synthesiser.py, llm\_client.py, main.py.
These can be copied directly from the source files.}

\section{Appendix: Sample LLM Outputs}

\placeholder{Insert raw JSON output from Stage 1 and Stage 2 for one representative
company (suggested: Bed Bath \& Beyond as it has both qualified opinion and covenant breach)
to illustrate the pipeline's intermediate outputs.}

\section{Appendix: Annual Report Sources}

All annual reports used in this study are publicly available from the following sources:

\begin{itemize}[leftmargin=1.5cm]
  \item \textbf{SEC EDGAR:} \href{https://www.sec.gov/cgi-bin/browse-edgar}{\texttt{https://www.sec.gov/cgi-bin/browse-edgar}}
  \item \textbf{HKEx披露易:} \href{https://www.hkexnews.hk}{\texttt{https://www.hkexnews.hk}}
  \item \textbf{HKICPA Auditor Examples:} \href{https://www.hkicpa.org.hk/professionaltechnical/assurance/example_auditors/SME-FRS-Jul2009.pdf}{\texttt{https://www.hkicpa.org.hk/.../SME-FRS-Jul2009.pdf}}
\end{itemize}

\end{document}
